\documentclass[11pt,a4paper]{article}
\usepackage[margin=19mm]{geometry}
\usepackage{amsmath,amssymb,booktabs,graphicx}
\usepackage[T1]{fontenc}
\usepackage{lmodern,microtype}
\usepackage[hidelinks]{hyperref}
\setlength{\parindent}{0pt}
\setlength{\parskip}{5pt}
\setlength{\abovedisplayskip}{6pt}
\setlength{\belowdisplayskip}{6pt}
\newcommand{\E}{\mathbb E}
\newcommand{\relu}{\operatorname{ReLU}}
\newcommand{\reportfigure}[3]{%
\par\addvspace{6pt}\noindent
\begin{minipage}{\linewidth}
\centering\includegraphics[width=#1\linewidth]{#2}\par
\vspace{3pt}
\refstepcounter{figure}
{\raggedright\fontsize{9}{11}\selectfont\setlength{\parskip}{0pt}%
\textbf{Figure \thefigure.} #3\par}
\end{minipage}\par\addvspace{6pt}}
\begin{document}
{\Large\bfseries Mitigating Neural Competition Through Overparameterization}\par

\subsection*{1. Do wider networks change the distribution of gradient agreement?}
Following Pier's suggestion, I compare minibatch-gradient product distributions
across widths, including their negative tails. The scalar-regression student and
half-MSE are
\[
f_\theta(x)=\frac1{\sqrt m}\sum_{r=1}^m a_r\relu(w_r^\top x+b_r),
\qquad L_Q(\theta)=\tfrac12\E_Q[(f_\theta-y)^2].
\]
Here $Q$ denotes the training data, a probe batch or the population, and
$\langle h,k\rangle_Q=\E_Q[hk]$. The first two studies use Gaussian inputs and an
eight-neuron teacher. Both layers and biases train, starting with zero readouts.
Width 8 already suffices to represent the target.

For independently sampled batches $B,C$ at fixed $\theta$, let $g_B=\nabla L_B$.
Then $\E_{B,C}[g_B^\top g_C]=\|\nabla L_{\rm train}\|^2$, by independence and
unbiasedness. Negative products predict increased loss on $C$ after a sufficiently
small step along $-g_B$. Since their raw distribution also reflects gradient scale,
I compare $g_B^\top g_C/\E_B\|g_B\|^2$ as well.

\reportfigure{1}{images/distributions.png}{Gradient-product CDFs across widths at similar training loss.
Top: raw products; bottom: products normalized by mean squared batch-gradient norm.
Negative values indicate first-order disagreement; shading shows uncertainty across starts.}

Matching near MSE .2 compares fitting progress rather than training time and
includes narrow Adam runs that never reach lower errors. This practical checkpoint
is not a theoretical threshold; it does not match residual structure or optimizer history.

Width changes raw scale more consistently than normalized shape. Normalized curves
largely overlap; fresh probes support modestly lighter tails for SGD/momentum, but
no clear Adam trend or universal skewness shift. Sampled changes look gradual, with
no resolved phase transition. This does not support a general account in which
width helps by removing gradient disagreement.

Pier's distinction between cause and consequence remains central. Later narrow
SGD states can agree more while fitting worse: batches may share an \emph{unresolved}
residual, without learning the missing features. Momentum's candidate directions
show a clearer early width effect than raw gradients, supporting his emphasis on
optimizer history. Shortening high-rate SGD steps reduces harm, and lowering the
rate from the same state removes a plateau: finite-step instability is another
explanation in that regime. These controls motivate testing which feature updates
are obstructed, rather than treating agreement itself as the mechanism.

\newpage
\subsection*{2. Can individually useful neuron updates obstruct one another?}
To move from data disagreement to neuron interaction, fix one proposed optimizer
step $\Delta\theta$ and apply its neuron blocks separately and jointly. Each block
contains $(w_r,b_r,a_r)$; the optimizer state and all block proposals stay fixed.
For $e=f_\theta-y$ and $d_r=f_{\theta+\Delta\theta_r}-f_\theta$, additivity of the
network gives $f_{\theta+\Delta\theta_r+\Delta\theta_s}-f_\theta=d_r+d_s$.
Expanding half-MSE on the same source batch $B$ therefore gives
\[
A_r=\langle e,d_r\rangle_B+\tfrac12\|d_r\|_B^2,\quad
J_{rs}=\langle d_r,d_s\rangle_B,\quad
\Delta L_{r,s}=A_r+A_s+J_{rs}.
\]
These are finite-step identities, including ReLU gate crossings. Positive $J_{rs}$
can mean ordinary diminishing returns: both neurons correct the same error.
The stricter event $A_r<0$, $A_s<0$, $A_s+J_{rs}>0$ instead says both help alone,
but adding $s$ after $r$ increases loss. No gradient is recomputed between them.

\begin{center}
\small
\begin{tabular}{lrr}
\toprule
Early momentum, matched MSE .2 & Width 8 & Width 16\\
\midrule
Strict events / eligible ordered pairs & 1.245\% & 0.253\%\\
Strict events per source batch & 0.481 & 0.474\\
Summed marginal harm / full-step gain & 0.394\% & 0.035\%\\
\bottomrule
\end{tabular}
\end{center}
{\fontsize{9}{11}\selectfont\raggedright
\textbf{Table 1.} Cohort means. Eligible pairs have $A_r,A_s<0$; marginal harm sums
$A_s+J_{rs}$ over strict events. The last row divides mean harm by the full step's
mean loss reduction within each start, then averages across starts.\par}

Width reduces the frequency and summed harm, but hardly the event count because
there are more pairs. More importantly, the harm is small relative to useful
progress. These overlapping pair alternatives cannot be added up as a causal
loss penalty; the table establishes a local obstruction, not its responsibility
for the final fitting gap.

Could accumulation make interaction decisive? On a fixed reference set $Q$, define
\[
C_+(T)=\sum_{t<T}\sum_{r<s}[J_{rs,t}]_+,\qquad [z]_+=\max(z,0).
\]
Unlike the strict-event test, this includes all positive overlaps. The distinction
matters: reducing $C_+$ could remove redundancy, while increasing it could accompany
more useful learning.

\reportfigure{.88}{images/cumulative_cost.pdf}{Width $8\to16$ in two retained eight-feature teacher cohorts.
Arrows connect cohort means; both axes use the same fixed reference set within each cohort.
Wider models fit better in both, but positive pair cost moves in opposite directions.
The cohorts have different training protocols; this is not an optimizer ranking.}

Momentum initially supports the cost account, but much of its reduction disappears
when matching accumulated isolated helpful movement rather than training duration:
earlier convergence itself reduces subsequent interaction. Adam is a stronger
counterexample: width improves fitting while \emph{increasing} $C_+$.
Thus a monotone ``less pair cost, better learning'' explanation fails across these
regimes. This does not rule out specific harmful interactions; it shows that the
aggregate mixes them with useful, redundant and trajectory-dependent movement.
The next hypothesis must identify \emph{which necessary feature is suppressed},
not merely how much neurons interact.

\newpage
\subsection*{3. Does competition suppress necessary signed feature roles?}
Can overlapping target components give a needed feature the wrong signal? Set
$b_r=0$ and derive the dynamics under \emph{unit-rate gradient flow} on half-MSE:
$\dot a_r=-\partial_{a_r}L_Q$, $\dot w_r=-\nabla_{w_r}L_Q$. Both layers move at
equal rates, without momentum, preconditioning or regularization. Assume
$\rho_r=\|w_r\|>0$ and differentiate away from ReLU gate boundaries
(almost surely for Gaussian inputs).

Let $\phi_u(x)=\relu(u^\top x)$, $u_r=w_r/\rho_r$ and
$c_r=a_r\rho_r/\sqrt m$, so $f=\sum_r c_r\phi_{u_r}$. Define
\[
H_f(u)=\langle y-f,\phi_u\rangle_Q
      =\langle y,\phi_u\rangle_Q-\sum_s c_s K_Q(u,u_s),
\qquad K_Q(u,v)=\langle\phi_u,\phi_v\rangle_Q.
\]
$H_f$ is the negative loss derivative for adding signed feature mass at $u$.
Differentiation and ReLU homogeneity give
\[
\dot a_r=\frac{\rho_r}{\sqrt m}H_f(u_r),\qquad
\dot\rho_r=\frac{a_r}{\sqrt m}H_f(u_r),\qquad
\dot c_r=\frac{a_r^2+\rho_r^2}{m}H_f(u_r).
\]
Consequently $\operatorname{sign}(c_r)H_f(u_r)<0$ contracts $|c_r|$ when $c_r\ne0$.
This is a local amplitude law, not a frozen-direction model or a convergence claim.
For a teacher $y=\sum_i\beta_i\phi_{v_i}$ and $f\approx0$, the signal for its
correctly oriented, correctly signed role is
\[
\chi_i=\operatorname{sign}(\beta_i)\sum_j\beta_jK_Q(v_i,v_j).
\]
Opposite-sign coactive components can outweigh its positive self contribution:
a needed role contracts. However, learning peers can remove that offending residual
and \emph{release} it. Initial suppression does not establish harmful peer dynamics
throughout training.

The experiment uses eight Gaussian-input teacher features in four orthogonal pair
planes, $60^\circ$ apart within each pair. Coefficients $(1,-3)$ create opposing
demands; $(1,3)$ changes only their signs before scaling to unit training variance.
Bias-free widths 8 and 32 use ordinary random weights and small random readouts.
\emph{Training uses finite-step SGD, momentum and Adam}, with unchanged rates across
widths. Gradient flow motivates the hypothesis; it is not the exact experimental
dynamics, especially for momentum and Adam.

\reportfigure{.88}{images/signed_roles.pdf}{Left: target signs change the width advantage (circles: SGD;
squares: momentum; triangles: Adam). Right: opposing targets lead narrow models to lose positive roles.
The dashed quota is necessary for exact population fitting, not an approximate-error bound.}

Opposing targets leave narrow MSE at roughly 2--5\% of target variance, with much
smaller wide errors; same-sign targets are easier. Exact population fitting requires
four units of each sign: one unit cannot supply two distinct teacher hyperplane
kinks. Narrow runs often lose positive roles, even from balanced signs; width leaves
enough units to carry them. Redundancy against role loss is a candidate explanation.

But $\chi_i$ fails to locate the width-benefit boundary, and enough signs do not
ensure the right directions. Flipping signs also changes target means and scaling,
so this control does not isolate suppression. The next test asks whether changing
coefficient learning changes the inadequate features that narrow networks acquire.

\newpage
\subsection*{4. Does early readout adaptation change which features are learned?}
To test that question, keep the bias-free task and write its training predictions
as $\Phi(W)a$, where $W=(w_1,\ldots,w_m)$ and
$\Phi_{ir}=\relu(w_r^\top x_i)/\sqrt m$. With $n$ training examples, define
\[
a_*\in\arg\min_a\|y-\Phi a\|^2,\qquad
R(W)=\frac{\|y-\Phi a_*\|^2}{2n}.
\]
No full-rank assumption is needed: $a_*$ may be nonunique but $\Phi a_*$ is the
unique projection of $y$ onto the current feature span. The normal equations
$\Phi^\top(y-\Phi a_*)=0$ cancel the cross term, giving
\[
L(W,a)=R(W)+\frac{\|\Phi(a-a_*)\|^2}{2n}.
\]
Thus $R$ isolates the limitation of the current representation; the remaining term
is readout lag. These are empirical training quantities using labels and current
features, with no teacher directions. For $f_*=\Phi a_*$ and empirical inner products,
\[
H_f(u)=H_{f_*}(u)-\langle\Phi(a-a_*),\phi_u\rangle_{\rm train}.
\]
This connects the two objects: readout lag changes the residual field. Removing
it jointly refits the coefficients, potentially changing their signs and magnitudes.
This acts beyond suppression of a specific role: it also changes hidden-layer
gradients directly through their coefficient factors.

On the opposing task, periodic least-squares readout corrections are applied to
width 8 during either the first or the last tenth of training. Both layers otherwise
train normally. Equal correction counts control intervention budget; the early arm
has more time for subsequent feature adaptation. A terminal refit separately tests
whether an ordinary failed model already contains adequate features.

\reportfigure{.88}{images/readout_timing.pdf}{Early readout correction gives a lasting rescue; late correction is much
weaker. Dashed lines show ordinary width-32 error. All errors are measured on held-out data before
any terminal refit; corrections use training data only.}

A terminal refit barely rescues ordinary narrow models, whereas repeated early
corrections close about 94\% of the ordinary mean narrow--wide error gap. The
benefit persists after corrections stop and in longer runs. Taken together, these
results locate the problem in \emph{feature formation}, not merely inaccurate final
coefficients, and show that early coefficient adaptation can substantially alter it.
The percentage measures gap closure, not a fraction uniquely caused by competition.

The controls limit the causal interpretation. Retaining optimizer moments preserves
most of the correction benefit, but resetting moments alone also substantially
rescues Adam; faster readout learning helps as well. There is therefore no unique
identification of residual-field suppression as the cause, nor evidence that width
and readout correction act through the same mechanism. Extra least-squares work
also makes this an intervention study, not a comparison of training efficiency.

Finally, early $R$ and readout lag did not improve forecasts across settings, and
$R$ usually decreases even in failing runs. Failure need not mean that features
become worse: they may improve along an inadequate path. The intervention supports
coefficient--feature feedback as consequential in this task, while leaving open the
central predictive question: which residual-field trajectories lose necessary
features, and when does extra width prevent that loss?
\end{document}
